\documentclass[a4paper,11pt]{article}
\usepackage[utf8]{inputenc}
\usepackage[T1]{fontenc}
\usepackage[english]{babel}
\usepackage[left=2.5cm,right=2.5cm,top=2cm,bottom=2cm]{geometry}
\usepackage{hyperref}
\usepackage{xcolor}
\usepackage{microtype}
\usepackage{lmodern}

\definecolor{primary}{RGB}{63, 150, 58}
\hypersetup{colorlinks=true, urlcolor=primary}

\begin{document}
\noindent
\textbf{Loïc Aron MBASSI EWOLO} \\
ENSEA -- Double Diplôme Ingénieur \\
\href{mailto:loicaron.mbassiewolo@ensea.fr}{loicaron.mbassiewolo@ensea.fr} \quad (+33) 7 59 52 33 98 \\
\href{https://github.com/Nameless0l}{github.com/Nameless0l} \quad \href{https://mbassiloic.tech}{mbassiloic.tech}

\vspace{1em}
\noindent Cergy, le 22 mars 2026

\vspace{1em}
\noindent
\textbf{Objet : Candidature -- Big Data et IA} \\
\textbf{Référence : Offre d'alternance -- Schneider Electric}

\vspace{1.5em}
\noindent Madame, Monsieur,

\vspace{0.8em}
La transformation digitale du secteur industriel repose sur la capacité à extraire de l'intelligence de fluxes massifs de données hétérogènes en temps réel. Vous recherchez un ingénieur capable de concevoir des architectures big data et des modèles IA pour l'industrie. Mon stage en cours à INRIA Saclay (avril-août 2026) me forme à la gestion et l'analyse de données massives, tandis que mes projets IoT appliqués (Urban Waste Management) m'ont exposé à l'intégration de capteurs distribués, au traitement temps réel et à la prédiction. Ces expériences me positionnent pour contribuer rapidement à vos initiatives data-driven.

\vspace{0.8em}
À INRIA Saclay (équipe TRIBE), je conçois et déploie un pipeline complet d'analyse de données à grande échelle : ingestion hétérogène (données structurées et textuelles), prétraitement et nettoyage, extraction de features et évaluation statistique. J'utilise Python (pandas, scikit-learn), ainsi que des techniques NLP avancées (embeddings, transformers). Ce travail m'a exposé aux défis réels du big data : gestion mémoire, optimisation requêtes, validation statistique rigoureuse. Le projet Urban Waste Management (premier prix CITS Hackathon) démontre mon aptitude à appliquer le ML à des données IoT réelles : prédiction des besoins de collecte via séries temporelles, optimisation des routes via algorithmes combinatoires et analyse coût-bénéfice. J'ai manipulé des données météo, de localisation, démographiques, consolidées pour une prédiction robuste.

\vspace{0.8em}
Ma formation au sein du programme MILA (Institut Québécois d'IA, juin-septembre 2025) approfondit mon expertise en deep learning appliqué. J'ai notamment étudié le phénomène de grokking (généralisation tardive) sur architectures profondes, travail combinant expérimentation rigoureuse et analyse mathématique. Ces compétences en ML avancé, conjuguées à une maîtrise des outils big data modernes (pandas, TensorFlow, PyTorch) et à une sensibilité aux défis opérationnels, font de moi un candidat capable de transformer des données industrielles en valeur mesurable.

\vspace{0.8em}
Je suis disponible dès septembre 2026 pour une alternance de 12 mois. Persuadé que l'avenir de l'industrie réside dans l'analyse intelligente de données massives distribuées, je suis motivé à concrétiser cette vision chez Schneider Electric.

\vspace{1.5em}
\noindent Veuillez agréer, Madame, Monsieur, l'expression de mes salutations distinguées.

\vspace{2em}
\noindent \textbf{Loïc Aron MBASSI EWOLO}

\end{document}
